\documentclass[12pt]{article}
\usepackage{amsmath,amssymb}

\begin{document}
\section*{{2020 AMC 10B Problems/Problem 17}}
Source: AoPS Wiki

\subsection*{Solution}
Consider the $10$ people to be standing in a circle, where two people opposite of each other form a diameter of the circle. Let us use casework on the number of pairs that form a diameter of the circle. Case 1: $0$ diameters There are $2$ ways: either $1$ pairs with $2$ , $3$ pairs with $4$ , and so on or $10$ pairs with $1$ , $2$ pairs with $3$ , etc. So $\boxed{\textbf{2}}$ cases. Case 2: $1$ diameter There are $5$ possible diameters to draw (everyone else pairs with the person next to them). So $\boxed{\textbf{5}}$ cases. Note that there cannot be $2$ diameters since there would be one person on either side that will not have a pair adjacent to them. The only scenario forced is when the two people on either side would be paired up across a diameter. Thus, a contradiction will arise. Case 3: $3$ diameters There are $5$ possible sets of $3$ diameters to draw.
Notice we are technically choosing the number of ways to choose a pair of two diameters that are neighbors to each other. This means we can choose the first diameter in the pair, and have only two diameters to choose from for the second in the pair. This means we have $5 \cdot 2=10$ possibilities for choosing 5 neighboring diameters. However, notice that there are duplicates, so we divide the $10$ possibilities by $2$ to get $\boxed{\textbf{5}}$ . Note that there cannot be a case with $4$ diameters because then there would have to be $5$ diameters for the two remaining people as they have to be connected with a diameter. A contradiction arises. Case 4: $5$ diameters There is only $\boxed{\textbf{1}}$ way to do this. Thus, in total there are $2+5+5+1=\boxed{\textbf{(C) }13}$ possible ways. ~Pearl2008 + Alzwang (Minor Edits)

\end{document}
